\originalchapter{8.333作业6：量子气体与低能激发}
\psource{Mehran Kardar，8.333 Fall 2013，PS6，截止2013年12月6日。选做题5、9--13也收入。}
\begin{exercise}[原题1(a)--(d)]\label{p:333-6-1}
原题提供代表表（He3能量栏本身漏印系数，保留缺口）：
\begin{center}\begin{tabular}{lllll}\toprule 粒子&$n$/m$^{-3}$&$m$/kg&$E_F$/eV&$T_F$/K\\\midrule 电子&$10^{29}$&$9\times10^{-31}$&4.4&$5\times10^4$\\ 核子&$10^{44}$&$1.6\times10^{-27}$&$10^8$&$1.1\times10^{12}$\\ 液He3&$2.6\times10^{28}$&$4.6\times10^{-27}$&$\text{缺系数}\times10^{-3}$&$10^1$\\\bottomrule\end{tabular}\end{center}
另给He3原子体积46.2 Å$^3$，与代表密度未严格一致；本题所求估算不使用缺印数值。
估算：(a)典型金属室温电子与声子热容比（表中电子$T_F\simeq5\times10^4$ K）；(b)室温中子热波长与晶体声子最短波长；(c)室温、常压H$_2$、He、O$_2$的$n\lambda^3$及出现量子效应的温度；(d)He4在1 K以下每质量热容$20.4T^3$ J kg$^{-1}$ K$^{-1}$，仅一支无简并低能激发，求$E(k)$。
\end{exercise}
\begin{solution}
(a)取每原子一个导电电子，电子$C_e=\pi^2N\kb T/(2T_F)$，高温晶格$C_{\rm ph}\simeq3N\kb$，故比$\pi^2T/(6T_F)\simeq0.0099$。若声子尚未到Dulong--Petit区，应使用实际Debye温度，不能当普适金属常数。

(b)$\lambda=h/\sqrt{2\pi m\kb T}$，300 K中子约1.00 Å；声子布里渊区边界$k\sim\pi/a$，最短波长约$2a$，若$a\sim2$ Å则约4 Å。比较尺度即可，不需要为“典型晶体”虚构晶格参数。(c)$n=P/(\kb T)$，取300 K、101325 Pa，三者依次$\lambda=(0.710,0.502,0.178)$ Å、$n\lambda^3=(8.76\times10^{-6},3.10\times10^{-6},1.37\times10^{-7})$。定密度冷却使$n\lambda^3\propto T^{-3/2}$，量子尺度依次$0.128,0.0638,0.00797$ K；定压则$T^{-5/2}$，依次$2.85,1.88,0.539$ K。实际气体可能先液化或冻结，这些是理想气外推。

(d)三维单支各向同性$E\propto k^s$使$C\propto T^{3/s}$，所以$s=1$。设$E=\hbar ck$，对Bose分布积分得$C/V=2\pi^2\kb^4T^3/(15\hbar^3c^3)$。还需质量密度；采用液He4约145 kg m$^{-3}$，$c=[2\pi^2\kb^4/(15\hbar^3\rho\,20.4)]^{1/3}\simeq240$ m/s。热容系数本身不足以无条件确定密度未知材料的声速；以上额外物性假设明确列出。
\end{solution}
\begin{exercise}[原题2(a)--(d)]\label{p:333-6-2}
太阳中心$T=1.6\times10^7$ K，氢质量密度$6\times10^4$、氦$10^5$ kg m$^{-3}$，完全电离。(a)电子、质子、$\alpha$粒子热波长；(b)理想近似下是否简并；(c)总气体压强；(d)辐射压强及主要支撑来源。
\end{exercise}
\begin{solution}
$n_p=\rho_H/m_p$，$n_\alpha=\rho_{He}/(4m_p)$，电中性$n_e=n_p+2n_\alpha$。热波长为$1.863\times10^{-11},4.349\times10^{-13},2.174\times10^{-13}$ m。按电子、质子自旋简并2，$n\lambda^3/g$为$0.213,1.48\times10^{-6},1.54\times10^{-7}$；核完全经典，电子弱简并，可有约百分之几修正，主阶仍可用经典压力。

$P_{\rm gas}=\kb T(n_e+n_p+n_\alpha)=2.58\times10^{16}$ Pa。$P_{\rm rad}=a_RT^4/3=1.65\times10^{13}$ Pa，$a_R=7.5657\times10^{-16}$ J m$^{-3}$ K$^{-4}$。气体约大1558倍，在给定中心条件主要物质压强支撑；这不取代整个恒星的静力结构方程。
\end{solution}
\begin{exercise}[原题3(a)--(e)]\label{p:333-6-3}
三维非相互作用同种粒子，自旋简并$g$，$\varepsilon=p^2/(2m)$、密度$n$。保留$n^2$项求(a)压力、(b)能量密度、(c)$C_V/N$、(d)$\alpha_P=V^{-1}(\partial_TV)_{P,N}$、(e)$C_P/N$，比较Bose/Fermi。
\end{exercise}
\begin{solution}
令$\eta=+1$为Bose、$-1$为Fermi，$f_s^\eta(z)=\sum_{l\ge1}\eta^{l-1}z^l/l^s$，$n=g\lambda^{-3}f_{3/2}^\eta$、$P=g\kb T\lambda^{-3}f_{5/2}^\eta$。消去$z$得到$B_2=-\eta\lambda^3/(g2^{5/2})$。记$x=nB_2$，$P=n\kb T(1+x)$、$U/V=3P/2$，$x\propto nT^{-3/2}$。定$V$微分给
\begin{equation}
C_V/N=\tfrac32\kb(1-x/2),\qquad
\alpha_P=T^{-1}(1-5x/2).
\end{equation}
后式由$\alpha_P=P_T/(nP_n)$且$P_T=n\kb(1-x/2)$、$P_n=\kb T(1+2x)$。$C_P-C_V=TV\alpha_P^2/\kappa_T=N\kb(1-3x)$，故$C_P/N=\tfrac52\kb(1-3x/2)$。Bose$x<0$提高这些首阶热容，Fermi相反；只在$|x|\ll1$成立。
\end{solution}
\begin{exercise}[原题4(a)--(c)]\label{p:333-6-4}
某中子星核候选相有两支Fermi激发$E_\pm(k)=\pm\hbar^2(|k|-k_F)^2/(2M)$，自旋$g=2$。令$q=|k|-k_F$，用$\dd^3k\simeq4\pi k_F^2\dd q$：(a)证明$\Delta U=(2gVk_F^2/\pi^2)\int_0^\infty E_+(q)[\e^{\beta E_+(q)}+1]^{-1}\dd q$；(b)闭式积分；(c)低温热容。
\end{exercise}
\begin{solution}
粒子与空穴各贡献同一能量，再由$q$正负两侧得到另一个2，合并$4\pi/(2\pi)^3$即题中系数。设$a=\hbar^2/(2M)$，$y=\beta aq^2$，积分变$\tfrac12 a^{-1/2}(\kb T)^{3/2}\int_0^\infty y^{1/2}/(\e^y+1)\dd y$。用Fermi积分，
\begin{equation}
\Delta U=\frac{gVk_F^2}{\pi^2\sqrt a}\Gamma(3/2)(1-2^{-1/2})\zeta(3/2)(\kb T)^{3/2}.
\end{equation}
$C_V$等于系数乘$\tfrac32\kb(\kb T)^{1/2}$。二次零能壳的态密度$\rho(E)\propto E^{-1/2}$使$C\propto\sqrt T$，区别普通光滑Fermi面$C\propto T$；壳近似要求热展宽$q\ll k_F$。
\end{solution}
\begin{exercise}[原题5(a)--(f)，选做]\label{p:333-6-5}
单层石墨烯$E_\pm=\pm t_iak$，双层$E_\pm=\pm(t_i^2a^2/2t_\perp)k^2$，$g=4$，面积$A_s$，未掺杂。(a)求$\mu(T)$；(b)证明粒子空穴能量积分$\Delta U=2gA_s\int\dd^2k/(2\pi)^2\,E_+(k)/(\e^{\beta E_+}+1)$；(c),(d)双层能与热容；(e)面积各半的混合样品双层电子热容何时主导；(f)声速约$2\times10^4$ m/s的声子与单层电子低温热容比较。
\end{exercise}
\begin{solution}
粒子空穴对称固定半填充，$f(E)+f(-E)=1$，故$\mu=0$；有限带宽对称也成立。升温正能电子与负能空穴两项相同给(b)。记$a_2=t_i^2a^2/(2t_\perp)$、$b=t_ia$，径向积分为
\begin{align}
\Delta U_2/A_s&=\frac{g}{2\pi a_2}\int_0^\infty\frac{E\dd E}{\e^{\beta E}+1}=\frac{g\pi(\kb T)^2}{24a_2},&C_2/A_s&=\frac{g\pi\kb^2T}{12a_2},\\
\Delta U_1/A_s&=\frac{g}{\pi b^2}\int_0^\infty\frac{E^2\dd E}{\e^{\beta E}+1}=\frac{3g\zeta(3)(\kb T)^3}{2\pi b^2},&C_1/A_s&=\frac{9g\zeta(3)\kb^3T^2}{2\pi b^2}.
\end{align}
相等条件$\kb T_\times=\pi^2t_\perp/[27\zeta(3)]\simeq0.304t_\perp$，更低温双层主导，且须低能抛物近似仍有效。二维线性色散声子也给$T^2$，但系数$\propto\sum_s c_s^{-2}$；$v_e=t_ia/\hbar$通常远大声速，声子系数大。自由膜弯曲支$\omega\propto k^2$还给$T$，张力可改变它。故不宜仅凭电子态密度就判定石墨烯总热容由电子主导。
\end{solution}
\begin{exercise}[原题6(a)--(e)]\label{p:333-6-6}
单粒子态密度$\rho(\epsilon)$、基态0、$N=\int_0^\infty\rho(\epsilon)/(\e^{\beta(\epsilon-\mu)}-\eta)\dd\epsilon$。(a)能量；(b)Bose临界温隐式式；(c)Fermi低温$\mu-E_F$；(d)能量增量；(e)热容。可用Sommerfeld式$\int gf=\int_0^\mu g+\pi^2g'(\mu)/(6\beta^2)+\cdots$。
\end{exercise}
\begin{solution}
$U=\int\epsilon\rho(\epsilon)f_\eta\dd\epsilon$。Bose激发容量$N=\int_{0^+}^\infty\rho(\epsilon)/(\e^{\epsilon/(\kb T_c)}-1)\dd\epsilon$，要在低能收敛，且基态宏观占据单独处理。

Fermi固定$N$先展开数积分，$\rho_F\delta\mu+\pi^2(\kb T)^2\rho'_F/6=0$。能积分含$E_F\rho_F\delta\mu$，恰好抵消其导数项中的$E_F\rho'_F$，因此
\begin{equation}
\delta\mu=-\frac{\pi^2(\kb T)^2}{6}\frac{\rho'_F}{\rho_F},\quad
\Delta U=\frac{\pi^2}{6}\rho_F(\kb T)^2,\quad C_V=\frac{\pi^2}{3}\rho_F\kb^2T.
\end{equation}
须$\rho_F$有限且在$\kb T$窗口平滑；前题二次零壳不满足，不能套用。
\end{solution}
\begin{exercise}[原题7(a)--(e)]\label{p:333-6-7}
无自旋量子点粒子的弱吸引把每态能降低$uN^2/(2V)$，$Z=Z_0\exp[\beta uN^2/(2V)]$、$u>0$。(a)关联压力；(b)证明$P_n=\kb T f_{3/2}^\eta/f_{1/2}^\eta-un$，$n\lambda^3=f_{3/2}^\eta$；(c)稀薄不稳定临界$u_c$首阶量子修正；(d)Fermi零温；(e)Bose降低温度时的行为。
\end{exercise}
\begin{solution}
$F=F_0-uN^2/(2V)$，$P=-F_V=P_0-un^2/2$。理想气$\partial_zP_0=(\kb T/\lambda^3z)f_{3/2}$、$\partial_zn=(\lambda^{-3}/z)f_{1/2}$，即(b)。机械稳定$P_n>0$，故$u_c=(\kb T/n)f_{3/2}/f_{1/2}$。低密度消去$z$得
\begin{equation}
u_c=\frac{\kb T}{n}[1-\eta n\lambda^3/2^{3/2}+O((n\lambda^3)^2)].
\end{equation}
Fermi零温$P_0=2nE_F/5$、$E_F\propto n^{2/3}$，故$u_c=2E_F/(3n)$；Pauli压力抗塌缩。Bose接近$z=1$时$f_{1/2}\to\infty$而$f_{3/2}$有限，$u_c\to0$。凝聚相理想压力独立$n$，任意吸引使负压缩率。该点吸引模型无短程排斥，不能由不稳定条件直接证明稳定液气共存，可能整体塌缩。
\end{solution}
\begin{exercise}[原题8(a)--(f)]\label{p:333-6-8}
$d$维各向异性谐陷阱，$U=m\sum_\alpha\omega_\alpha^2x_\alpha^2/2$，$\hbar\omega_\alpha\ll\kb T$，忽略零点与离散性。(a)累积态数及态密度；(b)巨正则粒子数；(c)能量；(d)高温每粒子能；(e)零温$\mu$；(f)低温热容尺度。按每个自旋组分计。
\end{exercise}
\begin{solution}
相空间椭球或振子占据数单纯形$\sum\hbar\omega_\alpha n_\alpha<E$的体积给$\mathcal N(E)=E^d/[d!\prod\hbar\omega_\alpha]$、$\rho(E)=E^{d-1}/[(d-1)!\prod\hbar\omega_\alpha]$。积分得到
\begin{equation}
N=f_d^-(z)\prod_\alpha\frac{\kb T}{\hbar\omega_\alpha},\quad U=d\kb T f_{d+1}^-(z)\prod_\alpha\frac{\kb T}{\hbar\omega_\alpha}.
\end{equation}
经典$z\ll1$比值1，$U/N=d\kb T$，每维动能和势能各$\kb T/2$。零温$E_F=[d!N\prod\hbar\omega_\alpha]^{1/d}$。代$\rho(E_F)=dN/E_F$入前题，$C_V=\pi^2dN\kb^2T/(3E_F)$，且需$\hbar\omega\ll\kb T\ll E_F$才能同时使用连续谱与低温展开。
\end{solution}
\begin{exercise}[原题9(a)--(d)，选做]\label{p:333-6-9}
密度$\rho_t$的光晶格陷位，每位空或最多一个局域粒子，能$-w<0$，对游离理想Bose谱$p^2/(2m)$无影响，总密度$n$。(a)陷位占据密度；(b)游离密度（自旋$g$）；(c)凝聚临界总密度并画小$w$情形；(d)凝聚相能量密度。
\end{exercise}
\begin{solution}
单陷位$\Xi_t=1+\e^{\beta(\mu+w)}$，故$n_t=\rho_t/[1+\e^{-\beta(\mu+w)}]$；硬占据约束使公式类似Fermi，与游离Bose统计不冲突。$n_g=g\lambda^{-3}\operatorname{Li}_{3/2}(z)$。凝聚$\mu\to0$给
\begin{equation}
n_*(T)=\frac{\rho_t}{1+\e^{-w/(\kb T)}}+g\lambda^{-3}\zeta(3/2),\quad
U/V=-wn_t(0,T)+\tfrac32g\kb T\lambda^{-3}\zeta(5/2).
\end{equation}
基态游离凝聚粒子能0。小$w$时低温陷位几乎满$n_*(0)=\rho_t$，升温陷位逐渐释放可使临界密度先降，再由游离激发容量$T^{3/2}$增长，允许再入凝聚。图取$g\lambda(T_\rho)^{-3}\zeta(3/2)=\rho_t$、$w/(\kb T_\rho)=0.01$。
\begin{center}\begin{tikzpicture}\begin{axis}[width=10cm,height=4.5cm,xlabel={$T/T_\rho$},ylabel={$n_*/\rho_t$},xmin=0,xmax=.8,ymin=.45,ymax=1.3]
\addplot[thick,structurecolor,domain=.0005:.8,samples=300]{1/(1+exp(-.01/x))+x^1.5};\end{axis}\end{tikzpicture}\end{center}
\end{solution}
\begin{exercise}[原题10(a)--(d)，选做]\label{p:333-6-10}
经典吸引的Mayer函数$f(r)=\exp(-\pi r^2/\lambda^2)$，傅里叶$f(\omega)=\lambda^3\exp[-\lambda^2\omega^2/(4\pi)]$。(a)环积分$R_l=\int\prod_i\dd^3q_i/V\prod_if(q_i-q_{i+1})$；(b)证明仅保留环时$\log Z=\log Z_0+V[\operatorname{Li}_{5/2}(x)-x+(1-2^{-5/2})x^2]/(2\lambda^3)$，$x=n\lambda^3$；(c)压力；(d)压缩率预示的凝聚不稳定。
\end{exercise}
\begin{solution}
(a)一个平移积分为$V$，其余卷积令各傅里叶动量相同，$R_l=V^{1-l}\int\dd^3\omega/(2\pi)^3[f(\omega)]^l=\lambda^{3l}/[l^{3/2}\lambda^3V^{l-1}]$。

(b)$l\ge3$环的顶点循环和反转对称数$2l$给$N^lR_l/(2l)$；$l=2$只有一条连边，精确项$N^2\lambda^3/(2V)$，不能把双边环重复计。合成题示多对数；$l=1$扣去不存在的自环。令$c_0=1-2^{-5/2}$，定$N$取体积导数得
\begin{equation}
P=n\kb T+\frac{\kb T}{2\lambda^3}[\operatorname{Li}_{5/2}(x)-\operatorname{Li}_{3/2}(x)-c_0x^2],\quad
\frac{P_n}{\kb T}=1-\frac{\operatorname{Li}_{1/2}(x)-\operatorname{Li}_{3/2}(x)}{2x}-c_0x.
\end{equation}
$x\to1^-$时$\operatorname{Li}_{1/2}$发散，因此必在收敛区间出现负压缩率；数值零点约0.70436。原题称凝聚转变，这里严格能推出环近似的机械不稳定；纯吸引无稳定核的原始经典模型可能塌缩，环求和本身不能证明有限密度稳定液相。
\end{solution}
\begin{exercise}[原题11(a)--(g)，选做]\label{p:333-6-11}
$d$维均匀无自旋Bose气$\epsilon=p^2/(2m)$。(a)巨势、密度用$f_s^+(z)=\operatorname{Li}_s(z)$表示；(b)$PV/U$；(c)$T_c(n)$；(d)凝聚以下热容；(e)全温热容图；(f)最大热容与经典极限比及$d=3$值；(g)$d\to2$与适用维数。
\end{exercise}
\begin{solution}
令$s=d/2$。对$-\log(1-z\e^{-\beta\epsilon})$作动量积分或逐项积分，$\mathcal G=-\kb TV\lambda^{-d}\operatorname{Li}_{s+1}(z)$、$n=\lambda^{-d}\operatorname{Li}_s(z)$，$U=sPV$，比$2/d$与经典相同。

只在$d>2$时$\zeta(s)$有限，$T_c=(2\pi\hbar^2/m\kb)[n/\zeta(s)]^{1/s}$。以下$z=1$，$C/(N\kb)=s(s+1)\zeta(s+1)/\zeta(s)(T/T_c)^s$。以上固定$N$求$z_T$得到
\begin{equation}
\frac C{N\kb}=s(s+1)\frac{\operatorname{Li}_{s+1}}{\operatorname{Li}_s}-s^2\frac{\operatorname{Li}_s}{\operatorname{Li}_{s-1}},\qquad
(T/T_c)^s=\zeta(s)/\operatorname{Li}_s(z).
\end{equation}
$d=3$的峰在$T_c$，$C_{\max}/C_\infty=(5/2)\zeta(5/2)/\zeta(3/2)=1.28378$。下面曲线直接用上式消去$z$，不是手工连接。
\begin{center}\begin{tikzpicture}\begin{axis}[width=10cm,height=4.5cm,xlabel={$T/T_c$},ylabel={$C/(N\kb)$},xmin=0,xmax=4,ymin=0,ymax=2.1]
\addplot[thick,structurecolor] coordinates {(0,0) (0.0144927536,0.00335976272) (0.0289855072,0.009502844) (0.0434782609,0.0174578392) (0.0579710145,0.0268781017) (0.0724637681,0.0375632891) (0.0869565217,0.0493782259) (0.101449275,0.0622236763) (0.115942029,0.076022752) (0.130434783,0.0907135933) (0.144927536,0.106245026) (0.15942029,0.122573795) (0.173913043,0.139662713) (0.188405797,0.157479358) (0.202898551,0.175995134) (0.217391304,0.195184576) (0.231884058,0.215024814) (0.246376812,0.235495161) (0.260869565,0.256576788) (0.275362319,0.278252457) (0.289855072,0.300506313) (0.304347826,0.323323706) (0.31884058,0.346691048) (0.333333333,0.370595687) (0.347826087,0.395025807) (0.362318841,0.41997034) (0.376811594,0.445418887) (0.391304348,0.471361658) (0.405797101,0.49778941) (0.420289855,0.524693402) (0.434782609,0.552065348) (0.449275362,0.579897381) (0.463768116,0.608182016) (0.47826087,0.636912124) (0.492753623,0.666080902) (0.507246377,0.69568185) (0.52173913,0.725708747) (0.536231884,0.756155634) (0.550724638,0.787016797) (0.565217391,0.818286746) (0.579710145,0.849960206) (0.594202899,0.882032101) (0.608695652,0.914497541) (0.623188406,0.947351811) (0.637681159,0.980590364) (0.652173913,1.01420881) (0.666666667,1.04820289) (0.68115942,1.08256852) (0.695652174,1.11730171) (0.710144928,1.15239861) (0.724637681,1.1878555) (0.739130435,1.22366875) (0.753623188,1.25983486) (0.768115942,1.29635041) (0.782608696,1.3332121) (0.797101449,1.37041669) (0.811594203,1.40796107) (0.826086957,1.44584218) (0.84057971,1.48405705) (0.855072464,1.5226028) (0.869565217,1.56147661) (0.884057971,1.60067573) (0.898550725,1.64019748) (0.913043478,1.68003926) (0.927536232,1.72019851) (0.942028986,1.76067274) (0.956521739,1.80145953) (0.971014493,1.84255649) (0.985507246,1.88396129) (1,1.92567168)};
\addplot[thick,structurecolor] coordinates {(1.0000009,1.92567097) (1.00000109,1.92567083) (1.00000132,1.92567065) (1.00000159,1.92567044) (1.00000193,1.92567018) (1.00000233,1.92566987) (1.00000281,1.92566949) (1.00000339,1.92566904) (1.0000041,1.92566849) (1.00000495,1.92566783) (1.00000598,1.92566703) (1.00000722,1.92566606) (1.00000873,1.92566489) (1.00001054,1.92566348) (1.00001273,1.92566178) (1.00001538,1.92565972) (1.00001858,1.92565724) (1.00002244,1.92565423) (1.00002711,1.92565061) (1.00003274,1.92564623) (1.00003955,1.92564094) (1.00004777,1.92563455) (1.0000577,1.92562683) (1.0000697,1.92561751) (1.00008419,1.92560625) (1.00010169,1.92559265) (1.00012284,1.92557622) (1.00014838,1.92555637) (1.00017923,1.9255324) (1.0002165,1.92550345) (1.00026152,1.92546849) (1.0003159,1.92542626) (1.0003816,1.92537525) (1.00046097,1.92531364) (1.00055684,1.92523924) (1.00067267,1.92514938) (1.00081261,1.92504086) (1.00098169,1.92490982) (1.00118597,1.92475157) (1.00143282,1.92456049) (1.00173111,1.92432977) (1.00209161,1.92405124) (1.00252733,1.923715) (1.00305404,1.92330915) (1.00369085,1.92281935) (1.0044609,1.92222836) (1.00539231,1.92151543) (1.0065192,1.9206556) (1.00788307,1.91961897) (1.00953444,1.91836965) (1.01153492,1.91686469) (1.01395981,1.91505283) (1.01690135,1.91287292) (1.02047285,1.91025233) (1.02481404,1.90710503) (1.03009795,1.90332955) (1.03654001,1.8988068) (1.04441011,1.89339788) (1.05404903,1.88694194) (1.0658912,1.87925439) (1.08049679,1.87012574) (1.09859845,1.85932138) (1.121171,1.84658293) (1.14953935,1.83163168) (1.18555187,1.81417478) (1.21293867,1.80191968) (1.24757113,1.78754758) (1.28012684,1.77506693) (1.31127106,1.7639655) (1.34142163,1.75392468) (1.3708629,1.74473057) (1.39980017,1.73623153) (1.42838874,1.72831553) (1.45675061,1.72089708) (1.48498475,1.71390925) (1.51317378,1.70729852) (1.54138833,1.70102131) (1.56969024,1.69504158) (1.59813471,1.68932916) (1.62677196,1.6838585) (1.65564845,1.67860773) (1.68480785,1.67355798) (1.71429182,1.66869284) (1.74414055,1.66399793) (1.77439334,1.65946058) (1.80508898,1.65506954) (1.83626614,1.6508148) (1.86796368,1.64668739) (1.900221,1.64267924) (1.93307826,1.63878304) (1.96657667,1.6349922) (2.00075876,1.63130071) (2.03566861,1.62770306) (2.07135217,1.62419424) (2.10785744,1.62076963) (2.14523485,1.61742498) (2.18353747,1.61415637) (2.22282138,1.61096019) (2.26314601,1.60783306) (2.3045745,1.60477186) (2.34717409,1.60177369) (2.39101662,1.59883583) (2.43617897,1.59595576) (2.48274366,1.5931311) (2.53079944,1.59035963) (2.58044197,1.58763928) (2.63177463,1.58496807) (2.68490939,1.58234417) (2.73996779,1.57976583) (2.7970821,1.57723141) (2.85639657,1.57473938) (2.91806897,1.57228824) (2.98227227,1.56987663) (3.04919661,1.56750322) (3.11905159,1.56516675) (3.19206894,1.56286606) (3.26850564,1.56059999) (3.34864756,1.55836749) (3.4328138,1.55616752) (3.52136185,1.55399912) (3.61469364,1.55186134) (3.71326295,1.54975331) (3.81758427,1.54767417) (3.92824358,1.5456231)};
\addplot[dashed,gray,domain=0:4]{1.5};\end{axis}\end{tikzpicture}\end{center}


一般$d$必须区分临界值与全温最大值。临界左值比$R_c=(s+1)\zeta(s+1)/\zeta(s)$在$s\downarrow1$趋0，但经典极限比恒为1，因此此时$R_c$显然不能是全温最大值。严格答案为上述两分支热容相对$sN\kb$的上确界；例如$s=1.05$时全温峰约1.00140，在$T>T_c$，不是临界值0.1594。随$d\to2^+$，$T_c\to0$、临界值趋0，而全温最大比趋1。$d\le2$均匀气在有限温无此BEC，必须直接用正常分支。原题若将三维“临界峰”结论外推所有$d$，会产生该矛盾。
\end{solution}
\begin{exercise}[原题12(a)--(d)，选做]\label{p:333-6-12}
三维非简并Bose气$P,T$与吸引二维表面接触，吸附每粒子降低能$u$。(a)气体化学势；(b)表面理想Bose面密度；(c)允许的最大压力$P_*$；(d)$\delta P=P_*-P\to0$的奇异性。
\end{exercise}
\begin{solution}
气相$\mu=\kb T\log[P\lambda^3/(\kb T)]$。面内最低能$-u$，有效逸度$z_s=\e^{\beta(\mu+u)}$，二维积分$n_2=\lambda^{-2}\operatorname{Li}_1(z_s)=-\lambda^{-2}\log(1-z_s)$。需$z_s<1$，所以$P_*=(\kb T/\lambda^3)\e^{-\beta u}$、$n_2\simeq\lambda^{-2}\log(P_*/\delta P)$。原题措辞“完全凝聚至表面”指与无限三维库平衡的极限：二维均匀理想气在有限温没有有限临界面密度，不能据此声称发生通常二维BEC。
\end{solution}
\begin{exercise}[原题13(a)--(d)，选做]\label{p:333-6-13}
平行板以$v$运动，超流激发谱$\epsilon(p)$。(a)由大质量板的动量能量改变得板系激发能$\epsilon-p\cdot v$；(b)激发净动量积分；(c)小$v$定义$P_\alpha=V\rho_nv_\alpha$并求$\rho_n$；(d)声子$\epsilon=cp$的低温贡献。
\end{exercise}
\begin{solution}
板给予$p$后自身动量减少$p$，能变$[(Mv-p)^2-(Mv)^2]/(2M)=-p\cdot v+O(M^{-1})$，总产生能即题式。$P=V\int\dd^3p/(2\pi\hbar)^3\,p\,n_B(\epsilon-p\cdot v)$；零速项奇函数0，线性项$-n_B'(\epsilon)p_\alpha p_\gamma v_\gamma$。角平均给
\begin{equation}
\rho_n=\frac13\int\frac{\dd^3p}{(2\pi\hbar)^3}p^2[-n_B'(\epsilon)]=\frac{2\pi^2(\kb T)^4}{45\hbar^3c^5}.
\end{equation}
最后以$x=\beta cp$积分，$\int_0^\infty x^4\e^x/(\e^x-1)^2\dd x=4!\zeta(4)$。这是质量密度，由$P/V$和$v$量纲验证，超流质量密度为总密度减$\rho_n$；激发数密度不能直接替代正常质量分数。
\end{solution}
